% GENERATED FILE. Edit canonical lessons, then run npm run generate:books.
% canonical-source-hash: fnv1a64:710d5c3b34c40fc1
% canonical-lessons: TA-C428-caral, TA-C428-turu, TA-C428-poli, TA-C428-tenku, TA-C428-vati

\chapter{A Spray of Rain}
\label{ch:ta-a-spray-of-rain}
\hlchaptermodality{428}

\begin{chapteropening}
\textbf{By the end of this chapter:} \emph{I can talk about drizzle and downpours, and water collecting and draining away.}

Complete the last lesson of chapter 428: I can talk about drizzle and downpours, and water collecting and draining away.
\end{chapteropening}

\section[\texorpdfstring{\ta{சாரல்} (cāral)}{சாரல் (cāral)}]{\ta{சாரல்} (cāral) — a spray of rain; drizzle}

\label{lesson:TA-C428-caral}



\begin{quote}
Before the new one: say the Tamil for to scorch, then the Tamil for to be sultry.
\end{quote}

\subsection*{You'll want to know: \ta{சாரல்}}
\textbf{\ta{சாரல்}} — \emph{cāral} — \textquotedblleft{}a spray of rain; drizzle\textquotedblright{}.

\begin{quote}
\ta{மலையில்} \ta{காலை} \ta{முழுவதும்} \ta{சாரல்} \ta{அடித்தது}.
\end{quote}

\begin{quote}
\emph{malaiyil kālai muḻuvatum cāral aṭittatu.}
\end{quote}

\begin{quote}
On the hill a fine rain blew all morning.
\end{quote}

\ta{சாரல்} \ta{அடி} is said of fine rain blown by the wind.

\begin{cousinweb}
\textbf{\ta{சாரல்}} (\emph{cāral}) is native Dravidian (DEDR 2482), \textquotedblleft{}drizzling rain\textquotedblright{}.
\end{cousinweb}

\begin{grammarlens}[title={What you've built}]
One more word for getting around.
\end{grammarlens}

\subsection*{Your turn}
\begin{itemize}
\raggedright
  \item \emph{Say it:} \emph{cāral}
  \item \emph{Say it:} \emph{cāral}, once more
  \item \emph{Say it:} \emph{malaiyil kālai muḻuvatum cāral aṭittatu.} — On the hill a fine rain blew all morning.
  \item \emph{Recall:} say the Tamil for to scorch, then the Tamil for to be sultry, then say \emph{cāral} again
\end{itemize}

\begin{tcolorbox}[breakable,colback=teal!4,colframe=teal!35!black,title={Before you move on}]
What does \ta{சாரல்} mean? (\textquotedblleft{}A spray of rain; drizzle\textquotedblright{}.) Say it once more.
\end{tcolorbox}
\section[\texorpdfstring{\ta{தூறு} (tūṟu)}{தூறு (tūṟu)}]{\ta{தூறு} (tūṟu) — to drizzle}

\label{lesson:TA-C428-turu}



\begin{quote}
Before the new one: say the Tamil for to be sultry, then the Tamil for a spray of rain.
\end{quote}

\subsection*{You'll want to know: \ta{தூறு}}
\textbf{\ta{தூறு}} — \emph{tūṟu} — \textquotedblleft{}to drizzle\textquotedblright{}.

\begin{quote}
\ta{வெளியே} \ta{மழை} \ta{தூறிக்கொண்டிருந்தது}.
\end{quote}

\begin{quote}
\emph{veḷiyē maḻai tūṟikkoṇṭiruntatu.}
\end{quote}

\begin{quote}
It was drizzling outside.
\end{quote}

Past (it …): \ta{தூறியது}, \emph{tūṟiyatu}.

The noun is \ta{தூறல்}, a drizzle.

\begin{cousinweb}
\textbf{\ta{தூறு}} (\emph{tūṟu}) is native Dravidian (DEDR 3398), \textquotedblleft{}to drizzle\textquotedblright{}; the same article gives \ta{தூறல்}, \textquotedblleft{}drizzling\textquotedblright{}.
\end{cousinweb}

\begin{grammarlens}[title={What you've built}]
One more word for getting around.
\end{grammarlens}

\subsection*{Your turn}
\begin{itemize}
\raggedright
  \item \emph{Say it:} \emph{tūṟu}
  \item \emph{Say it:} \emph{tūṟu}, once more
  \item \emph{Say it:} \emph{veḷiyē maḻai tūṟikkoṇṭiruntatu.} — It was drizzling outside.
  \item \emph{Recall:} say the Tamil for to be sultry, then the Tamil for a spray of rain, then say \emph{tūṟu} again
\end{itemize}

\begin{tcolorbox}[breakable,colback=teal!4,colframe=teal!35!black,title={Before you move on}]
What does \ta{தூறு} mean? (\textquotedblleft{}To drizzle\textquotedblright{}.) Say it once more.
\end{tcolorbox}
\section[\texorpdfstring{\ta{பொழி} (poḻi)}{பொழி (poḻi)}]{\ta{பொழி} (poḻi) — to pour down; to shower}

\label{lesson:TA-C428-poli}



\begin{quote}
Before the new one: say the Tamil for a spray of rain, then the Tamil for to drizzle.
\end{quote}

\subsection*{You'll want to know: \ta{பொழி}}
\textbf{\ta{பொழி}} — \emph{poḻi} — \textquotedblleft{}to pour down; to shower\textquotedblright{}.

\begin{quote}
\ta{இரவு} \ta{முழுவதும்} \ta{மழை} \ta{பொழிந்தது}.
\end{quote}

\begin{quote}
\emph{iravu muḻuvatum maḻai poḻintatu.}
\end{quote}

\begin{quote}
It poured all night.
\end{quote}

Past (it …): \ta{பொழிந்தது}, \emph{poḻintatu}.

Written and poetic; in speech, \ta{மழை} \ta{கொட்டியது}.

\begin{cousinweb}
\textbf{\ta{பொழி}} (\emph{poḻi}) is native Dravidian (DEDR 4556), \textquotedblleft{}to pour forth, to shower, as rain\textquotedblright{}.
\end{cousinweb}

\begin{grammarlens}[title={What you've built}]
One more word for getting around.
\end{grammarlens}

\subsection*{Your turn}
\begin{itemize}
\raggedright
  \item \emph{Say it:} \emph{poḻi}
  \item \emph{Say it:} \emph{poḻi}, once more
  \item \emph{Say it:} \emph{iravu muḻuvatum maḻai poḻintatu.} — It poured all night.
  \item \emph{Recall:} say the Tamil for a spray of rain, then the Tamil for to drizzle, then say \emph{poḻi} again
\end{itemize}

\begin{tcolorbox}[breakable,colback=teal!4,colframe=teal!35!black,title={Before you move on}]
What does \ta{பொழி} mean? (\textquotedblleft{}To pour down; to shower\textquotedblright{}.) Say it once more.
\end{tcolorbox}
\section[\texorpdfstring{\ta{தேங்கு} (tēṅku)}{தேங்கு (tēṅku)}]{\ta{தேங்கு} (tēṅku) — to stagnate; to collect (of water)}

\label{lesson:TA-C428-tenku}



\begin{quote}
Before the new one: say the Tamil for to drizzle, then the Tamil for to pour down.
\end{quote}

\subsection*{You'll want to know: \ta{தேங்கு}}
\textbf{\ta{தேங்கு}} — \emph{tēṅku} — \textquotedblleft{}to stagnate; to collect (of water)\textquotedblright{}.

\begin{quote}
\ta{தெருவில்} \ta{மழைநீர்} \ta{தேங்கியது}.
\end{quote}

\begin{quote}
\emph{teruvil maḻainīr tēṅkiyatu.}
\end{quote}

\begin{quote}
Rainwater collected in the street.
\end{quote}

Past (it …): \ta{தேங்கியது}, \emph{tēṅkiyatu}.

Of work too: \ta{வேலை} \ta{தேங்கியது}, the work piled up.

\begin{cousinweb}
\textbf{\ta{தேங்கு}} (\emph{tēṅku}) is native Dravidian (DEDR 3453), \textquotedblleft{}to fill, to rise to the brim, to collect\textquotedblright{}.
\end{cousinweb}

\begin{grammarlens}[title={What you've built}]
One more word for getting around.
\end{grammarlens}

\subsection*{Your turn}
\begin{itemize}
\raggedright
  \item \emph{Say it:} \emph{tēṅku}
  \item \emph{Say it:} \emph{tēṅku}, once more
  \item \emph{Say it:} \emph{teruvil maḻainīr tēṅkiyatu.} — Rainwater collected in the street.
  \item \emph{Recall:} say the Tamil for to drizzle, then the Tamil for to pour down, then say \emph{tēṅku} again
\end{itemize}

\begin{tcolorbox}[breakable,colback=teal!4,colframe=teal!35!black,title={Before you move on}]
What does \ta{தேங்கு} mean? (\textquotedblleft{}To stagnate; to collect (of water)\textquotedblright{}.) Say it once more.
\end{tcolorbox}
\section[\texorpdfstring{\ta{வடி} (vaṭi)}{வடி (vaṭi)}]{\ta{வடி} (vaṭi) — to drain away; to ebb}

\label{lesson:TA-C428-vati}



\begin{quote}
Before the new one: say the Tamil for to pour down, then the Tamil for to stagnate.
\end{quote}

\subsection*{You'll want to know: \ta{வடி}}
\textbf{\ta{வடி}} — \emph{vaṭi} — \textquotedblleft{}to drain away; to ebb\textquotedblright{}.

\begin{quote}
\ta{மழை} \ta{நின்றதும்} \ta{தண்ணீர்} \ta{வடிந்தது}.
\end{quote}

\begin{quote}
\emph{maḻai niṉṟatum taṇṇīr vaṭintatu.}
\end{quote}

\begin{quote}
As soon as the rain stopped, the water drained away.
\end{quote}

Past (it …): \ta{வடிந்தது}, \emph{vaṭintatu}.

With an object, to strain: \ta{சாதம்} \ta{வடி}, to drain rice.

\begin{cousinweb}
\textbf{\ta{வடி}} (\emph{vaṭi}) is native Dravidian (DEDR 5221), \textquotedblleft{}to drip, to trickle, to ebb\textquotedblright{}.
\end{cousinweb}

\begin{grammarlens}[title={What you've built}]
One more word for getting around.
\end{grammarlens}

\subsection*{Your turn}
\begin{itemize}
\raggedright
  \item \emph{Say it:} \emph{vaṭi}
  \item \emph{Say it:} \emph{vaṭi}, once more
  \item \emph{Say it:} \emph{maḻai niṉṟatum taṇṇīr vaṭintatu.} — As soon as the rain stopped, the water drained away.
  \item \emph{Recall:} say the Tamil for to pour down, then the Tamil for to stagnate, then say \emph{vaṭi} again
\end{itemize}

\begin{tcolorbox}[breakable,colback=teal!4,colframe=teal!35!black,title={Before you move on}]
What does \ta{வடி} mean? (\textquotedblleft{}To drain away; to ebb\textquotedblright{}.) Say it once more.
\end{tcolorbox}
